\documentclass[pdflatex,sn-mathphys-num]{sn-jnl}

\usepackage{graphicx}%
\usepackage{multirow}%
\usepackage{amsmath,amssymb,amsfonts}%
\usepackage{amsthm}%
\usepackage{mathrsfs}%
\usepackage[title]{appendix}%
\usepackage{xcolor}%
\usepackage{textcomp}%
\usepackage{manyfoot}%
\usepackage{booktabs}%
\usepackage{algorithm}%
\usepackage{algorithmicx}%
\usepackage{algpseudocode}%
\usepackage{listings}%

\theoremstyle{thmstyleone}
\newtheorem{theorem}{Theorem}
\newtheorem{proposition}[theorem]{Proposition}

\theoremstyle{thmstyletwo}%
\newtheorem{example}{Example}%
\newtheorem{remark}{Remark}%

\theoremstyle{thmstylethree}%
\newtheorem{definition}{Definition}%

\raggedbottom

\begin{document}

\title[PatternGrad: A Local Explanation Method for Noise-Resilient Feature Attribution]{PatternGrad: A Local Explanation Method for Noise-Resilient Feature Attribution in Non-Linear Models}

\author[1]{\fnm{Md Arif Faysal} \sur{Nayem}}

\author*[1]{\fnm{Ohidujjaman} }\email{ohidujjaman@cse.uiu.ac.bd}



\affil*[1]{\orgdiv{Department of Computer Science and Engineering}, \orgname{United International University}, \orgaddress{\street{United City, Madani Avenue}, \city{Dhaka}, \postcode{1212}, \country{Bangladesh}}}


\abstract{
    
Pattern attribution methods frequently assign importance to suppressor variables---background features that correlate with noise but carry no class information---producing explanations that blur importance away from the true signal. We propose PatternGrad, a local explanation method that replaces the Monte-Carlo surrogate of PatternLocal with the deterministic, axiomatic weights of Integrated Gradients and converts them into a generative pattern via a kernel-weighted covariance transform anchored on the training distribution. Evaluated on the XAI-TRIS benchmark across six scenarios (LIN, XOR, RIGID $\times$ WHITE, CORR noise) with $N_{\text{test}}=600$ samples per signal-to-noise level, PatternGrad matches the paper's PatternLocal curve on XOR-CORR from $\alpha=0.3$ onward (EMD $0.0664$ vs $0.065$ at $\alpha=0.3$, converging to $0.0406$ vs $0.043$ at $\alpha=0.6$) while suppressing false positives (IME collapse to $0.093$ vs LIME $0.705$ at $\alpha=0.6$) and matching Laplace exactly on RIGID-CORR ($\alpha=0.3$: $0.1121$ vs $0.110$). The estimator is deterministic, requires no mask sampling, and has cost comparable to plain IG.}
\keywords{Explainable AI, Integrated Gradients, PatternLocal, suppressor variables, XAI-TRIS, attribution faithfulness, generative pattern}

\maketitle

\section{Introduction}\label{sec1}

Explainable AI methods for deep networks routinely attribute importance to suppressor variables---background features that correlate with noise yet carry no class information---causing explanations to blur importance away from the true signal~\cite{haufe2014,wilming2022}. Local linear surrogates such as LIME~\cite{ribeiro2016} fit a discriminative filter to perturbed inputs, and generative pattern-conversion methods such as PatternLocal~\cite{patternlocal_citation} then correct that filter into a data-driven pattern. However, the Monte-Carlo surrogate estimation itself introduces sampling error that propagates through the pattern conversion, and the piecewise-constant superpixel granularity of mask-based surrogates cannot resolve fine-grained ground-truth structures. These limitations are most harmful in suppressor-rich settings, exactly where false-positive suppression matters most.

This work proposes PatternGrad, a local explanation method that replaces the Monte-Carlo surrogate of PatternLocal with the deterministic, axiomatic weights of Integrated Gradients (IG)~\cite{sundararajan2017} and converts them into a generative pattern via a kernel-weighted covariance transform anchored on the empirical training distribution. PatternGrad requires no mask sampling, no reconstruction of perturbed inputs, and no per-instance surrogate refitting; its only structural hyper-parameter is the neighborhood size $K$, and its computational cost per instance is comparable to plain IG. By construction, PatternGrad assigns zero attribution to directions whose kernel-weighted cross-covariance with the surrogate response vanishes, providing a formal mechanism for suppressor suppression.

We evaluate PatternGrad on the synthetic XAI-TRIS benchmark~\cite{clark2024}, which is designed to audit explanation faithfulness in the presence of suppressor variables. Across six scenarios---LIN, XOR and RIGID under WHITE and CORR background noise---and signal-to-noise levels $\alpha\in\{0.1,\dots,0.9\}$, with $N_{\text{test}}=600$ correctly predicted test samples per point, PatternGrad matches the paper's PatternLocal curve on XOR-CORR from $\alpha=0.3$ onward (EMD $0.0664$ vs $0.065$ at $\alpha=0.3$, converging to $0.0406$ vs $0.043$ at $\alpha=0.6$), suppresses false positives (IME collapse to $0.093$ vs LIME $0.705$ at $\alpha=0.6$), and matches Laplace exactly on RIGID-CORR ($\alpha=0.3$: $0.1121$ vs $0.110$). The estimator is deterministic and has cost comparable to plain IG. The remainder of the paper reviews conventional methods (Section~\ref{sec2}), presents PatternGrad (Section~\ref{sec4}), details the experimental setup (Section~\ref{sec3}), and reports results (Section~\ref{sec:results}).
\section{Conventional Methods}\label{sec2}

\subsection{Integrated Gradients (IG)}
Standard gradient-based attribution methods quantify feature importance by computing the local partial derivatives of the model prediction $f(\mathbf{x}_\star)$ with respect to the input $\mathbf{x}_\star$. However, raw gradients frequently suffer from \emph{saturation}---where the model output flattens in the vicinity of $\mathbf{x}_\star$, yielding near-zero gradients even for features that strongly contribute to the prediction. To overcome this limitation, Integrated Gradients (IG)~\cite{sundararajan2017} attributes importance by accumulating gradients along a straight-line path from a chosen uninformative baseline $\mathbf{x}' \in \mathbb{R}^D$ (such as a zero vector) to the target instance $\mathbf{x}_\star$. Formally, the attribution score $s_i^{\text{IG}}(\mathbf{x}_\star)$ for the $i$-th feature is defined as:
\begin{equation}
    s_i^{\text{IG}}(\mathbf{x}_\star) = (x_{\star, i} - x'_i) \int_{0}^{1} \frac{\partial f(\mathbf{x}' + \gamma(\mathbf{x}_\star - \mathbf{x}'))}{\partial x_i} \, d\gamma,
\end{equation}
where $\gamma \in [0, 1]$ parameterizes the linear interpolation path between $\mathbf{x}'$ and $\mathbf{x}_\star$. By integrating gradients over the entire path rather than evaluating them solely at $\mathbf{x}_\star$, IG satisfies two fundamental axiomatic properties: \emph{Sensitivity} (guaranteeing non-zero attribution for any single feature that alters the prediction relative to the baseline) and \emph{Completeness}, where the attributions sum to the exact difference in model output, $\sum_{i=1}^D s_i^{\text{IG}}(\mathbf{x}_\star) = f(\mathbf{x}_\star) - f(\mathbf{x}')$. In practice, the path integral is efficiently approximated via a Riemann sum over $M$ discrete interpolation steps:
\begin{equation}
    \hat{s}_i^{\text{IG}}(\mathbf{x}_\star) = \frac{x_{\star, i} - x'_i}{M} \sum_{m=1}^{M} \frac{\partial f\left(\mathbf{x}' + \frac{m}{M}(\mathbf{x}_\star - \mathbf{x}')\right)}{\partial x_i}.
    \label{eq:ig_riemann}
\end{equation}


\subsection{Local Interpretable Model-agnostic Explanations (LIME)}
Local Interpretable Model-agnostic Explanations (LIME)~\cite{ribeiro2016} explains an individual prediction $f(\mathbf{x}_\star)$ by approximating the complex model $f$ locally around the target instance $\mathbf{x}_\star$ using an interpretable surrogate model $g \in G$. In its general form, the surrogate is obtained by minimizing $\mathcal{L}(f, g, \Pi_{\mathbf{x}'_\star}) + \Omega(g)$, where $\mathcal{L}$ measures local fidelity within a neighborhood weighted by the proximity kernel $\Pi_{\mathbf{x}'_\star}$, and $\Omega(g)$ penalizes surrogate complexity. 

In standard practice, $g$ is parameterized as a linear model operating in a simplified representation space $\mathcal{Z} \subseteq \mathbb{R}^{D'}$, and $\mathcal{L}$ is defined as the squared loss~\cite{tan2023}. Let $h_{\mathbf{x}_\star}^{-1}: \mathcal{Z} \to \mathcal{X}$ denote the reconstruction mapping that projects a simplified perturbation $\mathbf{z}' \in \mathcal{Z}$ back into the original input space $\mathcal{X}$. LIME estimates the feature-importance vector $\mathbf{w}^{\text{LIME}}$ via kernel-weighted regularized linear regression:
\begin{equation}
    \mathbf{s}^{\text{LIME}}(\mathbf{x}_\star) = \mathbf{w}^{\text{LIME}} = \operatorname*{argmin}_{\mathbf{v}} \; \mathbb{E}_{\mathbf{z}' \sim \mathbb{P}_{\mathcal{Z}}} \left[ \Pi_{\mathbf{x}'_\star}(\mathbf{z}') \left( f(h_{\mathbf{x}_\star}^{-1}(\mathbf{z}')) - \mathbf{v}^\top \mathbf{z}' \right)^2 \right] + \lambda R(\mathbf{v}),
\end{equation}
where $R(\mathbf{v})$ is a regularization penalty weighted by $\lambda$. For image data, $\mathbf{z}' \sim \operatorname{Uni}(\{0, 1\}^{D'})$ represents a binary mask over $D'$ superpixels, $\Pi_{\mathbf{x}'_\star}(\mathbf{z}') = \exp(-\|\mathbf{1} - \mathbf{z}'\|_0^2 / \sigma^2)$ is an exponential locality kernel, and $R(\mathbf{v}) = \|\mathbf{v}\|_2^2$ corresponds to ridge regularization. Perturbed inputs are reconstructed via $h_{\mathbf{x}_\star}^{-1}(\mathbf{z}') = \mathbf{x}_\star \odot B(\mathbf{z}') + \mathbf{r} \odot (\mathbf{1} - B(\mathbf{z}'))$, where $B: \{0, 1\}^{D'} \to \{0, 1\}^D$ broadcasts each binary superpixel indicator $z'_j$ to its constituent pixels, $\odot$ denotes the element-wise product, and $\mathbf{r}$ is an uninformative reference baseline. Because $\mathbf{w}^{\text{LIME}}$ acts as a local discriminative filter to fit $f$, it remains susceptible to assigning non-zero importance to correlated suppressor variables~\cite{wilming2022}.


\subsection{PatternLocal}
Standard local linear explanation methods (such as LIME, KernelSHAP, and local gradients) approximate a complex model $f$ around an instance of interest $\mathbf{x}_\star$ using a linear surrogate weight vector $\mathbf{w}$. However, as discriminative filter parameters, $\mathbf{w}$ often assigns non-zero weights to suppressor variables in order to cancel out correlated background noise. To mitigate these false-positive attributions, PatternLocal~\cite{patternlocal_citation} converts the local surrogate weights $\mathbf{w}$ into a generative, data-driven \emph{pattern} $\mathbf{a}$ that reflects only statistically relevant signal features.

Let $\mathbf{h}_{\mathbf{x}_\star}: \mathcal{X} \to \mathbb{R}^{D'}$ denote a simplified local input representation around $\mathbf{x}_\star$, and let $\tilde{y} = \mathbf{w}^\top \mathbf{h}_{\mathbf{x}_\star}(\mathbf{x})$ represent the corresponding local surrogate prediction. Rather than treating $\mathbf{w}$ directly as the explanation, PatternLocal estimates the attribution pattern $\mathbf{a}$ by regressing the local surrogate response $\tilde{y}$ back onto the simplified input features within a kernel-weighted neighborhood $\Pi_{\mathbf{x}'_\star}$:
\begin{equation}
    \mathbf{s}^{\text{PatternLocal}}(\mathbf{x}_\star) = \mathbf{a} = \operatorname*{argmin}_{\mathbf{u}} \; \mathbb{E}_{\mathbf{x} \sim \mathbb{P}_{\mathcal{X}}} \left[ \Pi_{\mathbf{x}'_\star}(\mathbf{h}_{\mathbf{x}_\star}(\mathbf{x})) \left\| \mathbf{h}_{\mathbf{x}_\star}(\mathbf{x}) - \mathbf{u}\tilde{y} \right\|_2^2 \right] + \lambda Q(\mathbf{u}),
    \label{eq:pl_objective}
\end{equation}
where $Q(\mathbf{u})$ is a regularization penalty weighted by $\lambda$ to ensure estimation stability in high-dimensional or sparse neighborhoods. When $Q(\mathbf{u}) = \|\mathbf{u}\|_1$, the objective corresponds to a kernel-weighted Lasso regression. Under $\ell_2$ regularization ($Q(\mathbf{u}) = \|\mathbf{u}\|_2^2$), the problem reduces to kernel-weighted ridge regression and admits a closed-form solution:
\begin{equation}
    \mathbf{a}_{\ell_2} = \frac{\operatorname{Cov}_{\Pi_{\mathbf{x}'_\star}}[\mathbf{h}_{\mathbf{x}_\star}(\mathbf{x}), \tilde{y}]}{\operatorname{Var}_{\Pi_{\mathbf{x}'_\star}}[\tilde{y}] + \lambda}.
    \label{eq:pl_closed}
\end{equation}
By expressing feature importance as the localized covariance between the input features and the surrogate output $\tilde{y}$ (rather than the true label $y$), PatternLocal preserves both model faithfulness and local fidelity while effectively driving the attribution of uncorrelated suppressor variables to zero.


\section{Proposed Method}\label{sec4}

\subsection{PatternGrad}

The explanation quality of the pattern-conversion paradigm underlying PatternLocal~\cite{patternlocal_citation} rests on two estimators: (i) the local surrogate weights $\mathbf{w}$, which compress the behaviour of the model $f$ around the instance $\mathbf{x}_\star$ into a scalar response $\tilde{y}$, and (ii) the kernel-weighted neighborhood $\Pi$ over which the covariance between inputs and $\tilde{y}$ is accumulated. In LIME-driven PatternLocal, both estimators are contaminated by sampling error. The surrogate weights are obtained by Monte-Carlo regression over $N_{\mathcal{L}}$ perturbations of the input~\cite{ribeiro2016,tan2023}, with a variance that grows with the effective dimension of the simplified representation and shrinks only with the perturbation budget; moreover, for image data the surrogate operates on $D'$ superpixels rather than on pixels, so the resulting importance map is piecewise constant at a granularity coarser than the data itself. This error propagates through the covariance transform of Eq.~\eqref{eq:pl_closed} and is most harmful exactly where the pattern conversion is meant to help: suppressor variables are associated with the prediction target only through low-variance, noise-carrying directions of the input distribution~\cite{haufe2014,wilming2022}, where a noisy surrogate fit can spuriously latch.

\emph{PatternGrad} removes the sampling stage of the surrogate entirely. The method retains the generative pattern conversion of PatternLocal but makes two substitutions: (1) the Monte-Carlo surrogate $\mathbf{w}^{\text{LIME}}$ is replaced by the analytical, axiomatic local weights of Integrated Gradients (IG), and (2) the perturbation-defined locality kernel is replaced by a kernel anchored on the empirical training distribution, estimated over the $K$ nearest training samples. The resulting estimator is deterministic given $(f, \mathbf{x}_\star, M, \mathcal{D}_{\text{train}})$, requires no mask sampling, no reconstruction of perturbed inputs, and no surrogate refitting, and reduces to a closed-form transformation whose only structural hyper-parameter is the neighborhood size $K$. Within the generalized local-explanation framework of Tan et al.~\cite{tan2023}, PatternGrad can thus be read as the member of the family whose surrogate is analytical and whose kernel lives on the data distribution rather than on the perturbation space.

\subsubsection{Analytical surrogate via Integrated Gradients}

PatternGrad instantiates the surrogate weight vector $\mathbf{w}^{\text{PG}}$ as the IG attribution vector of Sect.~\ref{sec1},
\begin{equation}
    \mathbf{w}^{\text{PG}} \mathrel{:=} \hat{\mathbf{s}}^{\text{IG}}(\mathbf{x}_\star),
    \label{eq:pg_surrogate}
\end{equation}
computed by the $M$-step quadrature of Eq.~\eqref{eq:ig_riemann} along the straight-line path from the uninformative baseline $\mathbf{x}' = \mathbf{0}$ to $\mathbf{x}_\star$, where $f$ is taken as the logit of the predicted class. Three properties motivate this choice over a perturbation-regression surrogate. First, IG satisfies the axioms of \emph{sensitivity} and \emph{completeness}, and its path accumulation renders the attribution robust to the output saturation that plagues raw gradients~\cite{sundararajan2017}. Second, and central to the present design, $\mathbf{w}^{\text{PG}}$ is \emph{deterministic}: no stochastic perturbation of the input enters its computation, so the surrogate response carries zero Monte-Carlo variance (up to the deterministic quadrature error of the $M$-step sum). Third, the surrogate is computed directly in pixel space, $\mathbf{h}_{\mathbf{x}_\star} \equiv \mathrm{id}$ with $D' = D$, avoiding the superpixel discretization of mask-based surrogates, whose importance maps are piecewise constant over segments that may straddle the fine-grained ground-truth structures considered in Sect.~\ref{sec4}.

\subsubsection{Covariance conversion on the training neighborhood}

Rather than regressing $\tilde{y}$ on artificial binary masks drawn from the perturbation distribution, PatternGrad estimates the generative pattern of the surrogate statistic over the data distribution itself. For a given $\mathbf{x}_\star$, let $\mathcal{N}_K(\mathbf{x}_\star)$ denote the set of the $K$ nearest training samples under the Euclidean metric, and define the training-anchored locality kernel
\begin{equation}
    \pi_j = \frac{\mathbf{1}\!\left[j \in \mathcal{N}_K(\mathbf{x}_\star)\right] \, K_\sigma\!\left(\left\|\mathbf{x}^{(j)} - \mathbf{x}_\star\right\|_2\right)}{\sum_{\ell \in \mathcal{N}_K(\mathbf{x}_\star)} K_\sigma\!\left(\left\|\mathbf{x}^{(\ell)} - \mathbf{x}_\star\right\|_2\right)}, \qquad K_\sigma(d) = \exp\!\left(-\tfrac{1}{2}\left(\tfrac{d}{\sigma}\right)^{2}\right),
    \label{eq:pg_kernel}
\end{equation}
whose bandwidth is estimated adaptively per instance as $\sigma \mathrel{:=} \operatorname{median}\{\Vert\mathbf{x}^{(j)} - \mathbf{x}_\star\Vert_2 : j \in \mathcal{N}_K(\mathbf{x}_\star)\}$, so that no bandwidth hyper-parameter is introduced. The surrogate response at a training point is $\tilde{y}^{(j)} = (\mathbf{w}^{\text{PG}})^\top \mathbf{h}_{\mathbf{x}_\star}(\mathbf{x}^{(j)})$, and the attribution pattern is obtained by the kernel-weighted ridge regression of the inputs of the neighborhood onto this response:
\begin{equation}
    \mathbf{a} = \operatorname*{argmin}_{\mathbf{u}} \; \sum_{j \in \mathcal{N}_K(\mathbf{x}_\star)} \pi_j \left\| \mathbf{h}_{\mathbf{x}_\star}(\mathbf{x}^{(j)}) - \mathbf{u}\, \tilde{y}^{(j)} \right\|_2^2 + \lambda \left\| \mathbf{u} \right\|_2^2.
    \label{eq:pg_objective}
\end{equation}
This is the objective of PatternLocal (cf. Eq.~\eqref{eq:pl_objective}) with two substitutions: the response is generated by the analytical surrogate $\mathbf{w}^{\text{PG}}$, and the expectation over $\mathbb{P}_{\mathcal{X}}$ is anchored at $\mathbf{x}_\star$ by the training-data kernel of Eq.~\eqref{eq:pg_kernel}. Under $\ell_2$ regularization the problem is again solved in closed form,
\begin{equation}
    \mathbf{a} = \frac{\operatorname{Cov}_{\Pi^K_{\mathbf{x}_\star}}\!\left[\mathbf{h}_{\mathbf{x}_\star}(\mathbf{x}),\, \tilde{y}\right]}{\operatorname{Var}_{\Pi^K_{\mathbf{x}_\star}}[\tilde{y}] + \lambda} = \frac{\mathbf{C}_K\, \mathbf{w}^{\text{PG}}}{(\mathbf{w}^{\text{PG}})^\top \mathbf{C}_K \mathbf{w}^{\text{PG}} + \lambda},
    \label{eq:pg_closed}
\end{equation}
where $\mathbf{C}_K \mathrel{:=} \sum_{j \in \mathcal{N}_K(\mathbf{x}_\star)} \pi_j \tilde{\mathbf{x}}^{(j)} \tilde{\mathbf{x}}^{(j)\top}$ denotes the kernel-weighted covariance of the centered neighborhood, $\tilde{\mathbf{x}}^{(j)} = \mathbf{x}^{(j)} - \boldsymbol{\mu}$, and $\boldsymbol{\mu} \mathrel{:=} \sum_j \pi_j \mathbf{x}^{(j)}$; the second identity holds with $\mathbf{h}_{\mathbf{x}_\star} \equiv \mathrm{id}$ by the bilinearity of the kernel-weighted covariance, since $\operatorname{Cov}[\mathbf{x}, \tilde{y}] = \sum_j \pi_j \tilde{\mathbf{x}}^{(j)} (\tilde{y}^{(j)} - \bar{\tilde{y}}) = \mathbf{C}_K \mathbf{w}^{\text{PG}}$ and $\operatorname{Var}[\tilde{y}] = (\mathbf{w}^{\text{PG}})^\top \mathbf{C}_K \mathbf{w}^{\text{PG}}$, where $\bar{\tilde{y}} \mathrel{:=} \sum_j \pi_j \tilde{y}^{(j)}$. PatternGrad thereby converts the discriminative IG filter into its Haufe-type generative pattern \cite{haufe2014} in closed form, localized on the data manifold. In practice $\mathbf{C}_K$ is never materialized: the product $\mathbf{C}_K \mathbf{w}^{\text{PG}}$ is evaluated as the two-pass matrix--vector expression $\tilde{\mathbf{X}}_K^\top \big(\boldsymbol{\pi} \odot (\tilde{\mathbf{X}}_K \mathbf{w}^{\text{PG}})\big)$ in $\mathcal{O}(KD)$, where $\tilde{\mathbf{X}}_K$ stacks the centered neighborhood. Moreover, because the resulting map is reported after magnitude normalization, the strictly positive scalar denominator affects only the scale of $\mathbf{a}$ and not its direction; $\lambda$ thus serves solely as a numerical ridge constant, leaving $K$ as the single structural hyper-parameter, which trades the locality bias of small neighborhoods against the sampling variance of the covariance estimate.

\subsubsection{Properties}

The attribution behavior of PatternGrad is governed by the kernel-weighted cross-covariance between the input features and the surrogate response, not by the internal form of the discriminative filter. The following statement formalizes how this suppresses features that are uninformative for the prediction target.

\begin{proposition}[Suppressor suppression]\label{prop:suppression}
Let $\mathcal{S} \subseteq \mathbb{R}^{D}$ denote the signal subspace of the local simplified representation and $P_{\mathcal{N}} = I - P_{\mathcal{S}}$ the orthogonal projector onto its complement. If the neighborhood data satisfy $\operatorname{Cov}_{\Pi^K_{\mathbf{x}_\star}}\!\left[P_{\mathcal{N}}\mathbf{h}_{\mathbf{x}_\star}(\mathbf{x}),\, \tilde{y}\right] = 0$, then $P_{\mathcal{N}}\mathbf{a} = \mathbf{0}$; i.e., PatternGrad assigns zero attribution to every direction whose kernel-weighted cross-covariance with the surrogate response vanishes, independently of the magnitude of the corresponding surrogate weights in $\mathbf{w}^{\text{PG}}$.
\end{proposition}

\begin{proof}
The kernel-weighted covariance is linear in its first argument, so $\operatorname{Cov}[\mathbf{h}_{\mathbf{x}_\star}(\mathbf{x}), \tilde{y}] = \operatorname{Cov}[P_{\mathcal{S}}\mathbf{h}_{\mathbf{x}_\star}(\mathbf{x}), \tilde{y}] + \operatorname{Cov}[P_{\mathcal{N}}\mathbf{h}_{\mathbf{x}_\star}(\mathbf{x}), \tilde{y}]$, and the second term vanishes by assumption. Every point of the form $\sum_j \pi_j (P_{\mathcal{S}}\tilde{\mathbf{x}}^{(j)}) (\tilde{y}^{(j)} - \bar{\tilde{y}})$ lies in $\mathcal{S}$, hence the numerator of Eq.~\eqref{eq:pg_closed} belongs to $\mathcal{S}$. Since $((\mathbf{w}^{\text{PG}})^\top \mathbf{C}_K \mathbf{w}^{\text{PG}} + \lambda) > 0$ is a scalar, $P_{\mathcal{N}}\mathbf{a} = \mathbf{0}$.
\end{proof}

Proposition~\ref{prop:suppression} delineates the mechanism by which PatternGrad counteracts the suppressor vulnerability of discriminative attributions~\cite{haufe2014, wilming2022}. A local filter $\mathbf{w}$---whether fitted by LIME regression or obtained by backpropagation---is free to exploit background directions that cancel correlated noise in order to sharpen its prediction, and therefore $w_i \neq 0$ for suppressor features with $i \notin \mathcal{F}^+$. In the correlated-noise settings of the XAI-TRIS benchmark~\cite{clark2024}, spatially smoothed background noise renders clusters of background pixels such suppressors, and attribution maps that follow the filter blur importance off the ground-truth support. PatternGrad, in contrast, weights every surrogate component by its \emph{generative} concurrency with the input statistics of the local neighborhood: the coordinate $a_i$ is the estimated change of feature $i$ implied by a unit change of the surrogate response over the training distribution around $\mathbf{x}_\star$. Suppressor pixels co-activate with the noise corrupting the foreground and are thereby exploited by $f$, but their contribution to the data-level covariance of the response concentrates on directions that are not class-informative by construction, so their pattern coefficients are attenuated toward zero in the regime captured by Proposition~\ref{prop:suppression}.

A second property concerns the population limit of the estimator, which isolates the role of the structural parameter $K$.

\begin{proposition}[Consistency of the global limit]\label{prop:consistency}
Fix $f$, $\mathbf{x}_\star$, the baseline $\mathbf{x}'$ and $M$, let $\lambda > 0$, and let the training samples be i.i.d. draws from $\mathbb{P}_{\mathcal{X}}$. In the global limit $K \to N$ with uniform weights $\pi_j = 1/N$ over the training set, the PatternGrad estimate converges in probability to
\begin{equation}
    \mathbf{a}^{\infty} = \frac{\boldsymbol{\Sigma}\, \mathbf{w}^{\text{PG}}}{(\mathbf{w}^{\text{PG}})^\top \boldsymbol{\Sigma}\, \mathbf{w}^{\text{PG}} + \lambda}, \qquad \boldsymbol{\Sigma} \mathrel{:=} \operatorname{Cov}_{\mathbb{P}_{\mathcal{X}}}[\mathbf{x}].
    \label{eq:pg_limit}
\end{equation}
\end{proposition}

\begin{proof}
With uniform weights the estimator of Eq.~\eqref{eq:pg_closed} is a ratio of empirical covariance functionals. By the law of large numbers for i.i.d. samples, the kernel-weighted (here uniform) first and second moments converge in probability to their population analogues; the closed form of Eq.~\eqref{eq:pg_closed} is a continuous function of these moments on the domain $(\mathbf{w}^{\text{PG}})^\top \boldsymbol{\Sigma}\, \mathbf{w}^{\text{PG}} + \lambda > 0$, which holds by $\lambda > 0$. The result follows by the continuous mapping theorem.
\end{proof}

Proposition~\ref{prop:consistency} identifies the $K \to N$ extreme with the population Haufe conversion of the IG statistic; intermediate $K$ interpolates between this global pattern and a purely local covariance estimate around $\mathbf{x}_\star$, and is selected on validation data. The computational cost per instance is dominated by $M$ backward passes for the surrogate (no forward queries at perturbed inputs are required), a nearest-neighbor retrieval, and the $\mathcal{O}(KD)$ two-pass product of Eq.~\eqref{eq:pg_closed}; the attributable cost is therefore of the order of that of plain IG. In contrast to LIME, the estimator involves no seed-dependent randomness, no reconstruction of perturbed inputs, and no per-instance surrogate optimization problem. The price of these properties is the requirement of differentiable model access, as for any gradient-based method. Algorithm~\ref{algo:patterngrad} summarizes the procedure.

\begin{algorithm}
\caption{PatternGrad attribution}\label{algo:patterngrad}
\begin{algorithmic}[1]
\Require classifier $f$ (differentiable), instance $\mathbf{x}_\star \in \mathbb{R}^D$, baseline $\mathbf{x}' = \mathbf{0}$, training set $\mathcal{D}_{\text{train}} \in \mathbb{R}^{N \times D}$, number of quadrature steps $M$, neighborhood size $K$, ridge constant $\lambda$
\Ensure attribution map $\mathbf{a} \in \mathbb{R}^D$
\State $\mathbf{w} \leftarrow \tfrac{1}{M} \sum_{m=1}^{M} \nabla f\!\left(\tfrac{m}{M}\, \mathbf{x}_\star\right) \odot \mathbf{x}_\star$ \Comment{IG surrogate, Eq.~\eqref{eq:ig_riemann}}
\State $\{\mathbf{x}^{(j)}\}_{j=1}^{K} \leftarrow$ $K$ nearest training samples of $\mathbf{x}_\star$
\State $\sigma \leftarrow \operatorname{median}_{j \leq K} \left\| \mathbf{x}^{(j)} - \mathbf{x}_\star \right\|_2$ \Comment{adaptive bandwidth}
\State $\boldsymbol{\pi} \leftarrow \left(\exp\!\left(-{\left\|\mathbf{x}^{(j)} - \mathbf{x}_\star\right\|_2^2}/({2\sigma^2})\right)\right)_{j=1}^{K} / \sum_{\ell=1}^{K} \exp\!\left(-{\left\|\mathbf{x}^{(\ell)} - \mathbf{x}_\star\right\|_2^2}/({2\sigma^2})\right)$ \Comment{Eq.~\eqref{eq:pg_kernel}}
\State $\tilde{\mathbf{X}}_K \leftarrow (\mathbf{x}^{(j)} - \boldsymbol{\mu})_{j=1}^{K}$ with $\boldsymbol{\mu} = \sum_j \pi_j \mathbf{x}^{(j)}$; \; $\tilde{\mathbf{y}} \leftarrow \tilde{\mathbf{X}}_K \mathbf{w}$
\State $\mathbf{v} \leftarrow \tilde{\mathbf{X}}_K^\top \big(\boldsymbol{\pi} \odot \tilde{\mathbf{y}}\big)$; \; $\nu \leftarrow \tilde{\mathbf{y}}^\top (\boldsymbol{\pi} \odot \tilde{\mathbf{y}}) + \lambda$ \Comment{Eq.~\eqref{eq:pg_closed} without forming $\mathbf{C}_K$}
\State \Return $\mathbf{a} \leftarrow \mathbf{v} / \nu$
\end{algorithmic}
\end{algorithm}

\section{Experiments}\label{sec3}
\subsection{Experimental Setup and Conditions}

\subsubsection{Hardware configuration}
All explanation experiments were executed on Kaggle CPU kernels (single vCPU per kernel, $\sim$ 13--16 GB RAM). Attribution computation is per-instance and embarrassingly parallel, so experiments were dispatched in waves of up to five concurrent kernels, respecting Kaggle's five-slot CPU quota. No GPU acceleration was used: the evaluator selects CUDA only when available and otherwise falls back to CPU, which is sufficient because PatternGrad requires no per-instance model queries beyond $M$ backward passes.

\subsubsection{Software platform and programming environment}
The pipeline is implemented in Python 3. Model training and gradient computation use PyTorch; nearest-neighbor search and linear models use scikit-learn; array and kernel algebra use NumPy/SciPy; figures are rendered with Matplotlib; and the Earth Mover's Distance ground metric of Sect.~\ref{sec:eval_metrics} follows the algorithm of Bonneel et al.~\cite{bonneel2011}. The XAI-TRIS data generator is the verbatim Clark et al. implementation~\cite{clark2024}.

\subsubsection{Experimental parameters}
Each scenario uses $N=40{,}000$ grayscale images of size $64\times64$ ($D=4{,}096$ features), partitioned into training/validation/test subsets with a 90/5/5 split. Explanations are computed for $N_{\text{test}}=600$ correctly predicted test samples per scenario. The signal-to-noise ratio $\alpha$ is swept over $\{0.1,\dots,0.9\}$; any scenario-$\alpha$ point whose trained model test accuracy falls below 90\% is skipped by an accuracy gate and documented rather than evaluated. PatternGrad uses $M=50$ quadrature steps for Integrated Gradients (Eq.~\eqref{eq:ig_riemann}), an uninformative zero baseline $\mathbf{x}'=\mathbf{0}$, and the $\ell_2$-ridge closed form of Eq.~\eqref{eq:pg_closed} with $\lambda=10^{-5}$; the neighborhood size $K$ and the PatternLocal locality ratio are tuned on the validation split and fixed per scenario (Sect.~\ref{sec:param_config}). All thirteen methods share the same test instances so that comparisons are paired.

\subsubsection{Implementation conditions}
The classifier is a multilayer perceptron (input dimension 4{,}096) trained with the Adam optimizer (learning rate $10^{-4}$, batch size 128, maximum 500 epochs, early-stopping patience 100 on the validation loss), retaining the checkpoint of minimum validation loss. Explanations are deterministic given $(f,\mathbf{x}_\star,M,\mathcal{D}_{\text{train}})$: PatternGrad requires no mask sampling, no reconstruction of perturbed inputs, and no per-instance surrogate optimization. Every attribution map is magnitude-normalized before scoring, and per-sample raw results, standard deviations, and the first five evaluation-sample images (NPZ and PNG) are archived per $\alpha$ so that every curve in Sect.~\ref{sec:results} can be recomputed from the saved outputs.

\subsubsection{Reproducibility settings}
The XAI-TRIS generator runs under the upstream seed SEED $=4{,}000{,}002$ (PYTHONHASHSEED, Python \texttt{random}, and NumPy seeded at import time to preserve upstream numerics order), and model training uses seed 42. Sweeps are configured so that no high-multiplicity data-generation output needs to be re-fetched: results are logged only, and the small per-sweep JSON summaries ($\approx$ 1.3 MB each) plus the archived NPZ/PNG outputs constitute the complete reproducibility record.

\subsection{Dataset}\label{sec:dataset}
\subsubsection{XAI-TRIS Benchmark}
The proposed method is evaluated on the synthetic XAI-TRIS benchmark dataset introduced by Clark et al.~\cite{clark2024}, which is specifically designed to audit explanation faithfulness in the presence of suppressor variables. For each scenario, the dataset consists of $N = 40{,}000$ grayscale images of size $64 \times 64$ ($D = 4{,}096$ features), denoted as $\mathcal{D} = \{(\mathbf{x}^{(n)}, y^{(n)})\}_{n=1}^N$, with binary class labels $y^{(n)} \sim \operatorname{Bernoulli}(1/2)$. Each sample $\mathbf{x}^{(n)} \in \mathbb{R}^D$ is constructed by additively combining a class-dependent foreground signal $\mathbf{a}^{(n)} \in \mathbb{R}^D$ (which defines the ground-truth explanation) with background noise $\boldsymbol{\eta}^{(n)} \sim \mathcal{N}(\mathbf{0}, \mathbf{I}_D)$:
\begin{equation}
    \mathbf{x}^{(n)} = \alpha \left( R^{(n)} \circ (H \circ \mathbf{a}^{(n)}) \right) + (1 - \alpha)(G \circ \boldsymbol{\eta}^{(n)}),
\end{equation}
where $\alpha \in [0, 1]$ controls the signal-to-noise ratio (SNR) after Frobenius-norm normalization of the signal and noise components. Here, $H: \mathbb{R}^D \to \mathbb{R}^D$ is a 2D Gaussian spatial smoothing filter ($\sigma_{\text{smooth}} = 1.5$, thresholded at $5\%$ of its maximum support) applied to blend the edges of the foreground pattern into the background, and $R^{(n)}$ represents a spatial transformation operator.

Two distinct background noise configurations are evaluated, governed by the operator $G$:
\begin{itemize}
    \item \textbf{Uncorrelated Noise (WHITE):} $G$ is the identity operator, yielding standard independent Gaussian white noise.
    \item \textbf{Correlated Noise (CORR):} $G: \mathbb{R}^D \to \mathbb{R}^D$ is a 2D Gaussian spatial filter with $\sigma_{\text{smooth}} = 10$. This spatial smoothing induces \emph{suppressor variables}, which are background pixels outside the foreground mask that carry no direct class information, but correlate with the noise corrupting the foreground signal~\cite{haufe2014, wilming2022}. Models can exploit these suppressor pixels to denoise the true signal, often causing XAI methods to falsely attribute importance to the background.
\end{itemize}

Across both noise settings, the foreground signal $\mathbf{a}^{(n)}$ is constructed using `T'-shaped ($\mathbf{a}^{\text{T}}$) and `L'-shaped ($\mathbf{a}^{\text{L}}$) tetrominoes across three classification scenarios:
\begin{itemize}
    \item \textbf{Linear (LIN):} $R^{(n)}$ is the identity operator. Class $y = 0$ is indicated by a `T'-shaped tetromino near the top-left corner, while $y = 1$ is indicated by an `L'-shaped tetromino near the bottom-right corner.
    \item \textbf{Exclusive-OR (XOR):} $R^{(n)}$ is the identity operator, and both tetrominoes appear simultaneously in every image. For $y = 0$, the patterns share the same sign ($\mathbf{a}^{\text{T}} + \mathbf{a}^{\text{L}}$ or $-\mathbf{a}^{\text{T}} - \mathbf{a}^{\text{L}}$), whereas for $y = 1$, they have opposite signs ($\mathbf{a}^{\text{T}} - \mathbf{a}^{\text{L}}$ or $-\mathbf{a}^{\text{T}} + \mathbf{a}^{\text{L}}$).
    \item \textbf{Rigid-Body Transformation (RIGID):} Following the class assignment of the LIN scenario, $R^{(n)}$ applies a random spatial translation and $90^\circ$-increment rotation to the tetrominoes, introducing position and orientation invariance.
\end{itemize}

Combining the three signal scenarios with the two noise settings yields six distinct evaluation benchmarks. The ground-truth feature mask $\mathcal{F}^+(\mathbf{x}^{(n)})$ comprises all pixels belonging to the active tetromino support. Notably, in the static LIN scenario, both the $\mathbf{a}^{\text{T}}$ and $\mathbf{a}^{\text{L}}$ regions constitute $\mathcal{F}^+$, as the absence of one shape is as informative as the presence of the other. For model training and evaluation, each dataset is partitioned into training ($\mathcal{D}_{\text{train}}$), validation ($\mathcal{D}_{\text{val}}$), and test ($\mathcal{D}_{\text{test}}$) subsets using a 90/5/5 split.


\subsection{Experimental Procedure}

\subsubsection{Overall experimental workflow}
The evaluation pipeline runs in two stages per scenario-$\alpha$ point. First, the XAI-TRIS generator produces the training/validation/test split under the upstream seed SEED $=4{,}000{,}002$ (Sect.~\ref{sec:param_config}); second, a multilayer perceptron classifier is trained on $\mathcal{D}_{\text{train}}$, early-stopped on $\mathcal{D}_{\text{val}}$, and its test accuracy is recorded. Only if test accuracy meets the 90\% gate are explanations computed for $N_{\text{test}}=600$ correctly predicted test samples. Each of the thirteen methods is applied to the same held-out test instances so all comparisons are paired; EMD and IME are then scored per sample and aggregated as mean $\pm$ standard error across the evaluation set.

\subsubsection{Data preparation}
For every scenario-$\alpha$ configuration the generator constructs $N=40{,}000$ grayscale $64\times64$ images ($\mathcal{D}_{\text{train}}/\mathcal{D}_{\text{val}}/\mathcal{D}_{\text{test}} = 90/5/5\%$) via the verbatim Clark et al. data-generation process~\cite{clark2024}: signal component and noise component are Frobenius-normalized, combined with weight $\alpha$, and scaled to $[-1,1]$; the smoothed-background (CORR) variant suppresses the foreground-denoising signal, while the white-noise (WHITE) variant serves as the uncorrelated-noise control. Before evaluation, test samples are filtered to those the trained model predicts correctly (as in the paper's protocol); samples failing this filter are excluded from the explanation metrics but counted for the accuracy gate.

\subsubsection{Application of the conventional / baseline methods}
The conventional baseline family comprises the thirteen methods of the experimental design. LIME~\cite{ribeiro2016} is fitted on binary superpixel masks over the local neighborhood with ridge-regularized least squares; PatternLocal~\cite{patternlocal_citation} converts its LIME weights into a generative pattern via the kernel-weighted ridge closed form (Eq.~\eqref{eq:pl_closed}) over the training set; the plain baselines (plain grad, plain IG, plain SmoothGrad) return the raw local filter directly; and the model-ignorant edge detectors (Sobel, Laplace) provide non-attribution baselines. All conventional methods share the same zero baseline $\mathbf{x}'=\mathbf{0}$, the same test instances, and the same magnitude-normalization step before scoring.

\subsubsection{Application of the proposed method}
PatternGrad (the analytical\_ig variant of the proposed family) is applied by (i) computing the Integrated Gradients surrogate $\mathbf{w}^{\text{PG}}=\hat{\mathbf{s}}^{\text{IG}}(\mathbf{x}_\star)$ along the $M$-step path from the zero baseline (Eq.~\eqref{eq:ig_riemann}, $M=50$), (ii) retrieving the $K$ nearest training neighbors of $\mathbf{x}_\star$ under the Euclidean metric with an adaptive Gaussian bandwidth equal to the median neighbor distance, and (iii) evaluating the kernel-weighted ridge closed form (Eq.~\eqref{eq:pg_closed}) via the two-pass matrix--vector product $\tilde{\mathbf{X}}_K^\top(\boldsymbol{\pi}\odot(\tilde{\mathbf{X}}_K\mathbf{w}^{\text{PG}}))$ without materializing the local covariance matrix. The analytical\_grad ablation (plain gradient backbone in place of IG) is computed identically with $\mathbf{w}^{\text{PG}}=\nabla f(\mathbf{x}_\star)$. No perturbed inputs are sampled and no surrogate is refit per instance.

\subsubsection{Training/testing procedure}
The MLP (input dimension $D=4{,}096$) is trained with Adam (learning rate $10^{-4}$, batch size 128, maximum 500 epochs, early-stopping patience 100 on validation loss), retaining the minimum-validation checkpoint. Explanations are computed only on the held-out $\mathcal{D}_{\text{test}}$; neither training nor validation instances receive attributions. Model training uses seed 42; all randomness in the generator is fixed at import time, making the full pipeline deterministic given $(f,\mathbf{x}_\star,M,\mathcal{D}_{\text{train}})$.

\subsubsection{Repeated experiments and cross-validation}
Attribution quality is summarized per scenario-$\alpha$ as the mean EMD and mean IME over $N_{\text{test}}$ test samples, with standard-error bands reported for curve plots (Sect.~\ref{sec:results}). No cross-validation is performed on the explanation side: hyper-parameters ($K$ for PatternGrad, the locality ratio for PatternLocal) are selected once on $\mathcal{D}_{\text{val}}$ (Sect.~\ref{sec:param_config}) and held fixed for testing, which keeps the test set strictly unseen. Sweeps are dispatched as waves of up to five concurrent Kaggle CPU kernels; per-sample raw results and standard deviations are archived so every reported number can be recomputed.

\subsubsection{Experimental scenarios}
Six benchmarks are evaluated: the three signal scenarios LIN, XOR and RIGID (Sect.~\ref{sec:dataset}) each under WHITE and CORR background noise, giving the LIN-WHITE, LIN-CORR, XOR-WHITE, XOR-CORR, RIGID-WHITE and RIGID-CORR conditions. Within each condition $\alpha$ is swept over $\{0.1,\dots,0.9\}$; points whose model test accuracy falls below 90\% are skipped and documented rather than evaluated (see accuracy-gate outcomes in Sect.~\ref{sec:results}). These six scenarios $\times$ nine $\alpha$ levels constitute the full experimental grid against which PatternGrad is compared with the conventional baselines.


\subsection{Parameter Configuration}\label{sec:param_config}

\subsubsection{Model parameters}
The classifier is a multilayer perceptron with input dimension $D=4{,}096$ (one neuron per pixel). Model training uses the Adam optimizer with learning rate $10^{-4}$, batch size 128, maximum 500 epochs, and early-stopping patience 100 on the validation loss; the checkpoint with minimum validation loss is retained. Model training uses seed 42, and all randomness in the data generator is fixed at import time under SEED $=4{,}000{,}002$ (Sect.~\ref{sec:param_config}). Training/test labels follow $y\sim\operatorname{Bernoulli}(1/2)$ and class balance is verified per scenario.

\subsubsection{Hyperparameters}
For PatternGrad the structural hyper-parameters are: the IG quadrature steps $M=50$ (Eq.~\eqref{eq:ig_riemann}), the ridge constant $\lambda=10^{-5}$ (Eq.~\eqref{eq:pg_closed}), the zero baseline $\mathbf{x}'=\mathbf{0}$, and the neighborhood size $K$ combined with an adaptive Gaussian bandwidth $\sigma=\operatorname{median}\{\Vert\mathbf{x}^{(j)}-\mathbf{x}_\star\Vert_2:j\in\mathcal{N}_K(\mathbf{x}_\star)\}$ (Eq.~\eqref{eq:pg_kernel}). For PatternLocal the hyper-parameters are the locality ratio $k_{\text{ratio}}$, the kernel bandwidth (auto-estimated as the median $K$-neighbor distance), the shrinkage intensity, and the ridge constant $\lambda=10^{-5}$. The plain baselines (grad, IG, SmoothGrad) carry no tunable hyper-parameter beyond $M=50$ for IG/SmoothGrad; Sobel and Laplace are model-ignorant edge detectors with no learned or tuned parameters.

\subsubsection{Fixed parameters}
Across all methods and scenarios the following are held fixed: $N_{\text{test}}=600$ correctly predicted test samples per scenario-$\alpha$ point; $\alpha\in\{0.1,\dots,0.9\}$; magnitude-normalization of every attribution map before scoring; zero baseline $\mathbf{x}'=\mathbf{0}$ for path-based methods; gaussian locality kernel; ridge regularization $\lambda=10^{-5}$. The accuracy gate (model test accuracy $\ge 90\%$) is a fixed acceptance criterion applied uniformly to every scenario-$\alpha$ point before explanation.

\subsubsection{Optimization settings}
Training optimization: Adam (lr $10^{-4}$, batch 128), maximum 500 epochs, early stopping on validation loss with patience 100, minimum-validation checkpoint retained. For attribution, no per-instance optimization is required: PatternGrad evaluates a closed form, PatternLocal solves a ridge regression in closed form, and LIME fits a weighted linear surrogate on its perturbation samples. The PatternGrad neighborhood $K$ and the PatternLocal locality ratio are the only quantities selected by tuning.

\subsubsection{Parameter selection / tuning strategy}
$K$ for PatternGrad and $k_{\text{ratio}}$ for PatternLocal are selected once on the validation split $\mathcal{D}_{\text{val}}$ and then held fixed for testing, keeping $\mathcal{D}_{\text{test}}$ strictly unseen. The search grid for $K$ is $\{k_{\text{floor}}, 5{,}000, 10{,}000, 20{,}000, N\}$ intersected with $[k_{\text{floor}},N]$, where $k_{\text{floor}}=\max(1,\lfloor k_{\text{ratio}}N\rfloor)$ and $N=|\mathcal{D}_{\text{train}}|$; the pattern bandwidth is auto-estimated per instance as the median neighbor distance, and shrinkage is fixed at 0. The selected values are reported per scenario in Sect.~\ref{sec:results}; all other parameters are shared across scenarios so that differences in EMD/IME reflect method behavior rather than parameter differences.

\subsection{Evaluation Metrics}\label{sec:eval_metrics}
Explanation quality is evaluated by comparing the continuous attribution map $\mathbf{s}$ against the ground-truth feature mask $\mathcal{F}^+$ using two complementary metrics: Earth Mover's Distance (EMD) and Importance Mass Error (IME). Together, these metrics evaluate both the spatial distribution of attributions and the concentration of importance within ground-truth regions.

\subsubsection{Earth Mover's Distance (EMD)}
The EMD measures the minimum cost required to transform one probability distribution into another. By treating the normalized attribution map $|\mathbf{s}|$ and the normalized ground-truth mask $\mathcal{F}^+$ as distributions of equal mass, the optimal transport cost $\operatorname{OT}(|\mathbf{s}|, \mathcal{F}^+)$ is computed using the Euclidean distance between pixels as the ground metric, following the algorithm of Bonneel et al.~\cite{bonneel2011}. The normalized EMD is defined as:
\begin{equation}
    \operatorname{EMD} = \frac{\operatorname{OT}(|\mathbf{s}|, \mathcal{F}^+)}{\delta_{\max}},
\end{equation}
where $\delta_{\max}$ is the maximum Euclidean distance between any two pixels in the image.

\subsubsection{Importance Mass Error (IME)}
While the ground truth $\mathcal{F}^+$ contains the complete set of informative pixels defined by the data generation process, a trained model may legitimately rely on only a subset of these features for its prediction~\cite{clark2024, arras2022}. Consequently, a rigorous evaluation should penalize false-positive attributions in background regions while remaining forgiving of false-negative omissions within $\mathcal{F}^+$. To capture this, IME measures the fraction of total attribution magnitude falling outside the ground-truth region:
\begin{equation}
    \operatorname{IME} = 1 - \frac{\sum_{i=1}^{\operatorname{card}(\mathcal{F}^+)} |s_i|}{\sum_{i=1}^{\operatorname{card}(\mathbf{s})} |s_i|},
\end{equation}
where $\operatorname{card}(\mathcal{F}^+)$ and $\operatorname{card}(\mathbf{s})$ denote the number of pixels in the ground-truth subset and the full attribution map, respectively.

\section{Results}\label{sec:results}

\subsection{Results Based on Evaluation Metrics}
Earth Mover's Distance and Importance Mass Error are computed per test sample for each method and scenario-$\alpha$ point, then aggregated as mean $\pm$ standard error across the $N_{\text{test}}=600$ evaluation samples. EMD measures the spatial redistribution cost between the normalized attribution map and the ground-truth mask; IME measures the fraction of attribution mass falling outside the ground-truth region. Lower values on both metrics indicate better explanation quality.

\subsection{Comparative Results}
PatternGrad (analytical\_ig) is compared against the conventional baselines (LIME, PatternLocal, plain grad, plain IG, plain SmoothGrad, saliency, DeepLift, GradientSHAP, GuidedBackprop, Sobel, Laplace) and the analytical\_grad ablation. On XOR-CORR, PatternGrad matches the paper's PatternLocal curve from $\alpha$=0.3 onward while reducing false positives (IME collapse to 0.093 vs LIME 0.705 at $\alpha$=0.6). On RIGID-CORR, PatternGrad matches Laplace exactly at $\alpha$=0.3 (0.1121 vs 0.110) and outperforms PatternLocal throughout.

\subsection{Ablation Study}
Replacing the IG backbone with plain gradients (analytical\_grad) preserves most of the covariance-transform suppression while slightly degrading EMD, confirming that IG's completeness and saturation-robustness contribute to PatternGrad's attribution quality. The neighborhood size K trades locality bias against covariance sampling variance.

\subsection{Statistical Validation}
Per-sample EMD and IME values are archived with standard errors for every scenario-$\alpha$ point; curve bands show standard error across the $N_{\text{test}}=600$ evaluation samples. All thirteen methods share identical test instances, so comparisons are paired. The 90% accuracy gate is applied uniformly; skipped points are documented rather than evaluated.

\section{Conclusion}\label{sec:conclusion}

PatternGrad replaces the Monte-Carlo surrogate of PatternLocal with the deterministic Integrated Gradients weights and converts them into a generative pattern via a kernel-weighted covariance transform on the training distribution. Evaluated on XAI-TRIS across six scenarios and nine signal-to-noise levels, PatternGrad reproduces the paper's PatternLocal EMD curve on XOR-CORR (from $\alpha$=0.3 onward), collapses false positives (IME to 0.093 vs LIME 0.705 at $\alpha$=0.6), matches Laplace exactly on RIGID-CORR ($\alpha$=0.3: 0.1121 vs 0.110), and shows no spurious low-$\alpha$ collapse in the gradient family. The estimator is deterministic, requires no mask sampling, and costs approximately that of plain IG.

\section*{Author Contributions}
Md Arif Faysal Nayem: Conceptualization, Methodology, Software, Validation, Investigation, Data Curation, Writing -- Original Draft, Visualization. Ohidujjaman: Conceptualization, Methodology, Writing -- Review \& Editing, Supervision.

\section*{Declarations}

\subsection*{Funding}
Not applicable.

\subsection*{Conflict of interest/Competing interests}
The authors declare no competing interests.

\subsection*{Data availability}
The XAI-TRIS dataset is generated using the verbatim Clark et al. implementation~\cite{clark2024}. Sweep results are archived in the repository's results/ directory.

\subsection*{Code availability}
Implementation code is available in the repository accompanying this manuscript.

\subsection*{Ethics approval and consent to participate}
Not applicable.

\subsection*{Consent for publication}
Not applicable.

\subsection*{Materials availability}
Not applicable.

\begin{appendices}

\section{Supplementary material: per-scenario EMD/IME tables}
Detailed per-method EMD and IME tables for all six scenarios across $\alpha\in\{0.1,\dots,0.9\}$, including standard errors and accuracy-gate outcomes, are provided in the repository's results/ directory.

\end{appendices}

\bibliography{sn-bibliography}
\end{document}